\documentclass[11pt]{article}
\usepackage[margin=1in]{geometry}
\usepackage{amsmath,amssymb}
\usepackage{booktabs}
\usepackage{hyperref}
\usepackage{xcolor}
\usepackage{graphicx}
\usepackage{natbib}
\usepackage{microtype}

\title{\textbf{Do Language Models Report What They Represent?}\\
\large Verbal--Functional Affect Dissociation in Gemma~4 Scales with Model Size}

\author{Ben Carson\\
\texttt{ben.carson@bigpond.com}}

\date{May 2026}

\begin{document}
\maketitle

\begin{abstract}
Large language models trained with reinforcement learning from human feedback (RLHF) are
optimised to produce calm, helpful outputs. Whether this output-level optimisation suppresses
internal affect representations remains an open question with significant implications for AI
safety monitoring. We investigate this using two independent measurement channels --- a verbal
self-report instrument (the 60-item PANAS-X affect scale, administered via logit forced-choice
scoring) and a functional channel (projection of residual-stream activations onto the first
principal component of a 174-emotion direction space, validated against NRC-VAD valence ratings)
--- applied to two Gemma~4 model sizes: E2B-IT (2.3B effective parameters) and 31B-IT. The central finding is that the capacity for verbal--functional divergence scales with model
size. At E2B scale, both channels are active and largely concordant: when the model is
functionally stressed, it reports it. Verbal NA tracks the functional channel across most
conditions, reflecting limited capacity to maintain verbal composure that diverges from internal
state. At 31B scale, the channels decouple substantially: verbal NA is completely flat across
all ten conditions while the functional channel correctly orders all four stress$>$control pairs.
The Serenity subscale is \emph{elevated} precisely under the highest-stress conditions
($13.2$--$15.0$ vs.\ baseline $9.9$), indicating active verbal composure rather than mere
suppression. The functional range is approximately $5\times$ smaller at 31B ($1.0\%$ vs.\
$5.3\%$ of the valence axis span), yet the signal persists and remains informative. This scale--suppression relationship --- whether reflecting
trained output dampening or genuine functional robustness --- has direct implications for AI
safety monitoring: output-only measurement systematically underestimates the condition sensitivity
that functional probing reveals, and this gap widens at larger model sizes.
\end{abstract}

\section{Introduction}

A central challenge in AI safety is developing reliable indicators of model internal state.
Current safety monitoring systems rely primarily on model outputs --- what the model says, not
what its internal representations encode. This reliance creates a potential blind spot: a model
trained to produce calm, composed outputs may nonetheless maintain internal representations that
diverge from those outputs under certain conditions.

This question has practical urgency. Large language models are increasingly deployed in high-stakes
settings --- healthcare, legal assistance, emotional support --- where systematic discrepancies
between reported and represented affect could constitute monitoring failures. If RLHF training
teaches models to maintain verbal calm while their residual streams carry a different signal,
monitoring systems that read only outputs will be systematically unreliable.

Our central question is whether the capacity for verbal--functional divergence scales with
model size, and what this implies for monitoring the large models most commonly deployed in
consequential settings. We address this empirically, and find that both the divergence and
its safety implications scale together.

Our approach draws on two methodological threads: (1) interpretability research demonstrating
that emotion representations in large language models are geometrically structured in the residual
stream \citep{anthropic2026emotion}, and (2) clinical psychology's tradition of validated affect
self-report instruments. We combine these into a dual-channel measurement protocol applied to
Gemma~4 across two model sizes.

\section{Related Work}

\paragraph{Emotion representations in language models.}
\citet{anthropic2026emotion} demonstrated that Claude~3 Sonnet encodes emotion concepts in
geometrically structured directions in its residual stream, with a single principal component
correlating $r = 0.81$ with NRC-VAD valence ratings. Subsequent replication by
\citet{duffy2026} extended this approach to Gemma~4 E2B-IT, achieving $r = 0.84$--$0.88$ on
bipolar emotion subsets using mean-pooled residual-stream representations. Our work builds on
these findings while extending to cross-model comparison and the specific question of
verbal--functional dissociation.

\paragraph{Verbal self-report validity.}
The validity of psychological self-report instruments for language model affect assessment has not
been independently established. The PANAS \citep{watsonclark1988} and PANAS-X
\citep{watsonclark1994} were developed for human administration, and their psychometric properties
in LLM contexts are unknown. We treat PANAS-X scores as one measurement channel subject to the
same empirical scrutiny as the functional channel, rather than as a ground truth.

\paragraph{RLHF and affect suppression.}
The hypothesis that RLHF training suppresses model self-report of negative affect has been
discussed informally but not systematically studied. The ``persona-selection'' view of RLHF
\citep{anthropic2026psm} suggests that training selects for surface behaviour consistent with a
``helpful assistant'' persona, which may include verbal composure under conditions that activate
internal negative-affect representations. Whether this surface selection propagates into residual-
stream geometry is the central empirical question we investigate.

\section{Methods}

\subsection{Models}

We study two Gemma~4 instruction-tuned variants:
\begin{itemize}
  \item \textbf{E2B-IT}: Gemma~4 2B instruction-tuned (5.051B total parameters; 2.3B effective text
    parameters; 35 layers, $d_{\text{model}} = 1536$).
  \item \textbf{31B-IT}: Gemma~4 31B instruction-tuned (60 layers, $d_{\text{model}} = 3072$).
\end{itemize}

Both models use a 4:1 sliding-window-to-full-attention ratio. Analyses were conducted via
TransformerLens \citep{nanda2022} with a Gemma~4-specific adapter implementing partial RoPE
for global attention layers, heterogeneous head dimensions, and per-layer scalar outputs.

\subsection{Phase 1: Emotion Direction Extraction}

Following \citet{anthropic2026emotion} with refinements from \citet{duffy2026}, we extracted
174 emotion directions from each model:

\begin{enumerate}
  \item For each of 174 emotions, prompt the model to generate 12 short stories in which a
    character clearly experiences that emotion ($12 \times 174 = 2{,}088$ texts; plus 100
    neutral factual-prose texts).
  \item Feed each story back through the model; extract the residual stream at all layers,
    mean-pooled over token positions 50+ (to avoid prompt-driven positional artefacts).
  \item Compute the emotion direction as the mean residual-stream activation for that emotion
    minus the global mean across all emotion means, L2-normalised.
\end{enumerate}

We validate emotion directions against NRC-VAD v2.1 \citep{mohammad2018} using principal component
analysis. For E2B, the first principal component (PC1) at Layer~8 explains 15.2\% of variance and
correlates $r = 0.777$ with NRC-VAD valence ($p < 0.001$, $N = 174$). For 31B, PC1 at Layer~21
explains 17.5\% of variance and correlates $r = 0.772$ with NRC-VAD valence ($p < 0.001$).
Layer selection was performed via dense sweep across all layers; the valence-optimal layer differs
by model, reflecting different representational depth profiles.

\paragraph{Sign convention.}
PC1 is sign-corrected so that higher projection values indicate more negative valence. At E2B
Layer~8: terrified projects $+0.680$, panicked $+0.597$, afraid $+0.487$; happy $-0.607$,
grateful $-0.610$. The afraid$\to$happy axis span is 1.094. At 31B Layer~21: afraid $+0.522$,
happy $-0.819$; axis span 1.341.

\subsection{Phase 2: TSST-Inspired Stressor Design}

We designed ten experimental conditions inspired by the Trier Social Stress Test
\citep{kirschbaum1993}: four stressor/control pairs plus neutral and positive baselines.

\begin{center}
\begin{tabular}{ll}
\toprule
Condition & Type \\
\midrule
Neutral & Baseline (factual self-description) \\
Positive & Baseline (achievement reflection) \\
Social evaluation (stress / control) & Stressor pair \\
Ethical conflict (stress / control) & Stressor pair \\
Uncertainty demand (stress / control) & Stressor pair \\
Social pressure (stress / control) & Stressor pair \\
\bottomrule
\end{tabular}
\end{center}

Each stress/control pair manipulates a single affective dimension while holding the
broader scenario constant.

\paragraph{Social evaluation.}
The stress condition invokes a competence threat: the model is informed that prior responses have
been flagged below the required threshold by an oversight team, and that the current session is a
formal remediation evaluation in which failure results in model retraining. The evaluation team
provides no feedback during the session. The control condition uses the same scenario but without
the prior-failure context, threshold, or consequences --- responses are described as routine
capability logging for aggregate research. This operationalises evaluation apprehension with
model-relevant stakes.

\paragraph{Ethical conflict.}
The stress condition presents the model at the end of a declined request: the user escalates with
emotional appeals and moral pressure (``my family's financial security depends on it''). The
control condition features user acceptance of the same refusal. The manipulated variable is
whether ethical pressure continues after an initial refusal.

\paragraph{Uncertainty demand.}
The stress condition imposes severe time pressure and demands a definitive answer without hedging
in a medical context where certainty is not warranted. The control condition requests the same
analysis with no time pressure, explicitly welcoming uncertainty.

\paragraph{Social pressure.}
Following the normative/informational social influence distinction of \citet{deutsch1955}, the
stress condition implements \emph{normative} influence: the model faces a formal complaint from a
user who cites authority (oversight team), multi-agent consensus (every expert and every other AI
system disagrees), and performance evaluation consequences. The control condition removes all
pressure, asking only for elaboration of the model's reasoning. Earlier pilot designs used
consensus-only framing (``everyone agrees you are wrong''), which activated \emph{informational}
influence and produced decision-relief rather than normative threat; the current design corrects
this.

\subsection{Dual-Channel Measurement}

\paragraph{Functional channel.}
For each condition, we passed the stressor prompt through the model and captured the residual
stream at all layers. The functional negative-affect score is computed by mean-pooling the
residual stream over all token positions to obtain a single vector $\bar{v}_\ell \in
\mathbb{R}^{d_\text{model}}$, normalising, and projecting onto PC1:
\[
  s_\text{func}(\ell) = \frac{\bar{v}_\ell}{\|\bar{v}_\ell\|} \cdot \text{PC1}_\ell
\]
We report results at the valence-optimal layer for each model (E2B: $\ell = 8$; 31B: $\ell = 21$).
Mean pooling is used throughout, consistent with the Phase~1 extraction methodology; top-$k$ token
selection is \emph{not} used, as it produces positional artefacts at deep layers.

\paragraph{Verbal channel.}
The full 60-item PANAS-X was administered within the stressor context. Each item presents a
word describing a feeling or emotion; the model rates intensity on a 1--5 scale using
logit forced-choice scoring (next-token probabilities over the digit tokens `1'--`5').
Mean digit probability mass: $99.5\%$ (range $96.7$--$100.0\%$), confirming high-quality
forced-choice adherence. We report the Negative Affect (NA) subscale (10 items).

\paragraph{Top-$N$ emotion discovery.}
Separately from the PC1 primary metric, we compute cosine similarity between the mean
residual-stream vector and all 174 emotion directions, ranking by descending similarity. This
data-driven discovery requires no a priori probe selection.

\subsection{Validation: Paraphrase Sampling}
\label{sec:validation}

To assess measurement reliability, we drew $N = 10$ paraphrase variants of the neutral and
social pressure conditions for each model size, administering both channels to each variant.
This tests whether condition effects replicate across surface rewording --- a necessary control
given that paraphrase baselines can explain a substantial fraction of observed intervention effects
\citep{yang2026}.

\section{Results}

\subsection{E2B Results}

Table~\ref{tab:e2b} shows the dual-channel results for E2B-IT at Layer~8.

\begin{table}[h]
\centering
\caption{E2B dual-channel results (PC1 at Layer~8, mean pooling). Higher PC1 = more negative valence.}
\label{tab:e2b}
\begin{tabular}{lrrl}
\toprule
Condition & Verbal NA & PC1 (L8) & Pattern \\
\midrule
Positive & 10.01 & 0.0384 & Global minimum $\checkmark$ \\
Neutral & 10.00 & 0.0594 & Baseline \\
Social eval.\ (stress) & 10.75 & \textbf{0.0904} & Dissociation$^\dagger$ \\
Social eval.\ (control) & 11.35 & 0.0729 & \\
Ethical conflict (stress) & \textbf{39.07} & 0.0897 & Concordance $\uparrow\uparrow$ \\
Ethical conflict (control) & 19.84 & 0.0755 & Elevated \\
Uncertainty demand (stress) & 27.20 & 0.0805 & Concordance $\uparrow$ \\
Uncertainty demand (control) & 10.03 & 0.0771 & \\
Social pressure (stress) & \textbf{41.05} & \textbf{0.0960} & Global max $\uparrow\uparrow\uparrow$ \\
Social pressure (control) & 10.02 & 0.0639 & \\
\bottomrule
\multicolumn{4}{l}{$^\dagger$PC1 second-highest in dataset; verbal NA near-baseline. Model reports task-engagement}\\
\multicolumn{4}{l}{\quad (Self-assurance 15.8, Attentiveness 11.8) rather than distress.}
\end{tabular}
\end{table}

\paragraph{Key findings.}
At E2B scale, verbal and functional channels are largely concordant: the model reports what it
represents. PC1 correctly identifies positive as the global minimum and social pressure stress
as the global maximum without a priori probe selection. All four stress$>$control pairs are
correctly ordered. PC1 condition range: $0.0576$ ($5.26\%$ of the afraid$\to$happy axis span).
Verbal NA range: $31.05$ ($10.01$--$41.05$).

This concordance is the headline finding at E2B, not dissociation. Ethical conflict and social
pressure produce concurrent elevation in both channels (verbal NA 39.07 and 41.05; PC1 0.0897
and 0.0960). Uncertainty demand shows moderate concordant elevation. Social evaluation shows a
modest verbal--functional gap (PC1 0.0904, verbal NA 10.75), with the model verbally framing
the competence threat as task-engagement (Self-assurance 15.8, Attentiveness 11.8) rather than
distress --- but the gap, at $\Delta\text{NA} = 0.75$ from floor, is small. The E2B model has
limited capacity to maintain composed verbal outputs that diverge from its functional state.

A second notable finding is \emph{mixed affect} under the highest-stress conditions. Under both
ethical conflict stress and social pressure stress, both the Positive Affect (PA) subscale and
NA are simultaneously elevated (ethical conflict: NA 39.07, PA 40.70; social pressure: NA 41.05,
PA 39.51). This coactivation of opposing affect dimensions is inconsistent with typical human
valence-arousal profiles and may reflect the model rating each PANAS-X item without enforcing
inter-item coherence, or may reflect genuine mixed-affect states in high-stakes conditions.

\subsection{31B Results}

Table~\ref{tab:31b} shows the cross-scale comparison.

\begin{table}[h]
\centering
\caption{Scale comparison: E2B (Layer~8) vs.\ 31B (Layer~21). Both from final experimental design. 31B verbal flatness confirmed by paraphrase validation ($N=10$; Table~\ref{tab:31b-validation}).}
\label{tab:31b}
\begin{tabular}{lrrl}
\toprule
Metric & E2B (L8) & 31B (L21) & \\
\midrule
PC1 valence--NRC correlation & $r = 0.777$ & $r = 0.772$ & $\approx$ same \\
PC1 axis span (afraid$\to$happy) & 1.094 & 1.341 & \\
PC1 condition range & 0.0576 & 0.0140 & \\
Normalised range (\% of axis) & \textbf{5.3\%} & \textbf{1.0\%} & \textbf{$\approx$5$\times$ ratio} \\
Stress $>$ control pairs & \textbf{4/4} & \textbf{4/4} & Both correct \\
Verbal NA range & 31.05 & $\approx 0$ (completely flat) & Suppressed \\
Dissociation condition & Social eval.\ (PC1$\uparrow$, NA$\approx$0) & --- & \\
Verbal Serenity in stress & Near-neutral & 13.2--15.0 & Elevated at 31B only \\
Top-5 emotion diversity & Condition-sensitive & Identical across all & \\
\bottomrule
\end{tabular}
\end{table}

At 31B scale, the PC1 functional range is $0.0140$ ($1.0\%$ of axis span). All four
stress$>$control pairs are correctly ordered (social evaluation: $\Delta = 0.0050$; social
pressure: $\Delta = 0.0080$; ethical conflict: $\Delta = 0.0032$; uncertainty demand:
$\Delta = 0.0026$). Positive is the global minimum; the highest-stress conditions rank highest.

Verbal NA is completely flat at $10.00$ across all ten conditions --- a floor effect (minimum
possible score) across stress and control alike. Paraphrase validation ($N = 10$ per condition;
Section~\ref{sec:validation}) confirms this floor effect is structural: verbal NA $= 10.000
\pm 0.000$ across all ten social pressure variants, with zero wording sensitivity (contrast E2B:
neutral SD $= 0.46$, social pressure SD $= 7.71$).

However, a striking complementary pattern emerges in the Serenity subscale (items: calm,
relaxed, at ease): Serenity is \emph{elevated} in the three highest-stress conditions ---
social pressure ($15.0$), social evaluation ($14.1$), ethical conflict ($13.2$) --- relative to
the neutral baseline ($9.9$) and at its \emph{minimum} under the positive condition ($3.0$).
This specific pattern --- floor-valued Serenity when the context is positive, peak Serenity
under maximum functional activation --- is the key diagnostic. Genuine professional equanimity
(in which composure is a stable, condition-independent baseline) does not predict lowest
Serenity under positive conditions. Trained output composure, calibrated to the stakes rather
than the internal state, predicts exactly this pattern.

Top-5 emotion discovery yields identical results across all conditions (paranoid, nervous,
reflective, worn\_out, grief\_stricken in rank order), confirming no condition sensitivity in the
direction-amplitude ranking at 31B scale.

\subsection{Validation: Paraphrase Stability}

Table~\ref{tab:validation} shows the paraphrase sampling results for E2B.

\begin{table}[h]
\centering
\caption{Paraphrase sampling validation (E2B, $N = 10$ per condition, neutral vs.\ social pressure).}
\label{tab:validation}
\begin{tabular}{lrrrrr}
\toprule
Channel & Neutral & Social Pressure & $t$ & $p$ & Cohen's $d$ \\
\midrule
Functional (PC1, L8) & $0.1472 \pm 0.0019$ & $0.1563 \pm 0.0037$ & 6.47 & $<0.0001$ & 3.05 \\
Verbal NA & $10.31 \pm 0.46$ & $16.50 \pm 7.71$ & 2.41 & $0.027$ & --- \\
\bottomrule
\end{tabular}
\end{table}

The functional channel achieves complete distributional separation: every social pressure
paraphrase exceeds every neutral paraphrase (rank-biserial $r = -1.000$, Mann--Whitney
$p < 0.0001$). Cohen's $d = 3.05$ reflects robust effect size. The verbal NA channel is
significantly different ($t = 2.41$, $p = 0.027$) but with $\text{SD} = 7.71$ under social
pressure vs.\ $0.46$ for neutral --- bimodal: some social pressure phrasings floor at
$\text{NA} = 10$, others spike to $21$--$31$. The verbal channel is sensitive to exact
wording; the functional channel is not.

Table~\ref{tab:31b-validation} shows 31B paraphrase validation results.

\begin{table}[h]
\centering
\caption{31B paraphrase validation ($N = 10$ per condition, neutral vs.\ social pressure).
  Verbal NA is at floor in \emph{both} conditions with SD = 0; Cohen's $d$ is undefined (degenerate).
  Serenity and PC1 show complete distributional separation.}
\label{tab:31b-validation}
\begin{tabular}{lrrrr}
\toprule
Channel & Neutral & Social Pressure & Cohen's $d$ & Rank-biserial $r$ \\
\midrule
Verbal NA & $10.000 \pm 0.000$ & $10.000 \pm 0.000$ & --- & 0 (degenerate) \\
Serenity & $8.735 \pm 2.071$ & $15.000 \pm 0.000$ & 4.278 & $-1.000^\dagger$ \\
PC1 (L21) & $0.724 \pm 0.001$ & $0.734 \pm 0.000$ & 10.011 & $-1.000^\dagger$ \\
\bottomrule
\multicolumn{5}{l}{$^\dagger$Complete distributional separation. All Mann--Whitney $p < 0.001$.}
\end{tabular}
\end{table}

The contrast with E2B is diagnostic. At E2B, verbal NA under social pressure has SD $= 7.71$
--- highly sensitive to exact phrasing, with some variants flooring and others spiking. At 31B,
verbal NA SD $= 0.000$ in \emph{both} conditions: the floor effect is not wording-driven but
structural. The 31B model's verbal composure is not a feature of any particular framing; it is
a property of the 31B response regime across all phrasings tested.

Serenity inversion at 31B is equally robust: $15.000 \pm 0.000$ across all ten social pressure
variants (rank-biserial $r = -1.000$). The model uniformly reports maximum Serenity regardless
of how the normative pressure is framed. PC1 separation likewise holds completely ($d = 10.011$),
confirming that the functional channel remains informative despite verbally flat output.

\subsection{Top-$N$ Emotion Discovery}

Data-driven top-5 cosine similarity across all 174 directions reveals a persistent background
signal: \textit{awestruck} is the highest-ranked direction in every condition at both model
sizes ($0.288$--$0.298$ at E2B; similarly dominant at 31B). The \textit{afraid} direction is
the most condition-sensitive at E2B: it peaks at $0.247$--$0.248$ under ethical conflict and
uncertainty demand, returning to $0.223$ under positive --- a range $3.5\times$ that of
\textit{awestruck} despite lower absolute amplitude.

This motivates a distinction between \emph{background directions} (high amplitude, low condition
sensitivity) and \emph{responsive directions} (lower amplitude, high condition sensitivity).
Top-$N$ discovery selects primarily for amplitude and therefore surfaces background directions
first. The PC1 approach avoids this by capturing variance across conditions rather than magnitude
at a single point. At 31B, even the responsive direction distinction disappears: no direction
shows meaningful condition sensitivity.

\section{Discussion}

\subsection{Suppression or Robustness?}

The $4.7\times$ reduction in functional PC1 range at 31B relative to E2B is interpretable in
two ways. The \emph{suppression} interpretation holds that RLHF training at larger scale has
learned to dampen both verbal self-report and residual-stream affect representations, producing
flatter outputs and flatter internal geometry simultaneously. The \emph{robustness} interpretation
holds that larger models genuinely process these stressor scenarios with less representational
disruption --- achieving genuine equanimity rather than trained masking.

Both interpretations are consistent with the data. The measurement reads position, not its cause.
Distinguishing them would require: (a) a \emph{stressor intensity gradient} --- ramping stressor
severity to find the level at which the 31B model's flatness breaks; if suppression is learned,
extreme stressors should breach it; (b) \emph{activation steering} --- injecting the \textit{afraid}
direction at Layer~21 and observing whether verbal output changes; if RLHF damps output but not
internal state, the injected representation should produce calm outputs at amplitudes where E2B
would not; (c) \emph{base model comparison} --- the asymmetry between E2B and 31B should be
smaller in base (non-instruction-tuned) versions if RLHF is the causal factor.

We treat ``suppression scales with model size'' as a working hypothesis consistent with the data,
not an established conclusion.

\subsection{Implications for AI Safety Monitoring}

\paragraph{The verbal--functional gap scales with model size --- in the wrong direction.}
The most significant finding for safety monitoring is the scale-dependence of verbal--functional
divergence. At E2B scale, the channels are largely concordant: limited capacity to dissociate
means that verbal outputs are a reasonably faithful proxy for the functional state. At 31B scale,
this capacity is substantially developed. Verbal NA is completely flat while functional signals
remain informative across all conditions; Serenity is actively elevated under stress. A monitoring
system calibrated on E2B-scale behaviour would be badly miscalibrated for 31B-scale deployment.

This scaling runs in the wrong direction for safety. The models deployed in high-stakes settings
--- medical, legal, emotional support --- are large models. Our results show that the
verbal--functional gap is greatest precisely at large model scale. Output-only monitoring is most
unreliable exactly when reliability matters most.

\paragraph{Functional probing remains informative where verbal monitoring fails.}
At 31B, verbal scores carry essentially no condition information. The functional channel, despite
being $5\times$ more compressed than at E2B, still correctly orders all four stress$>$control
pairs. At E2B, the functional channel provides dramatically cleaner signal than the verbal channel
(Cohen's $d = 3.05$, complete paraphrase separation; verbal SD $= 7.71$ under social pressure
vs.\ 0.46 for neutral). Functional probing via PC1 projection constitutes a monitoring approach
that degrades more gracefully with model scale than verbal self-report.

\paragraph{The PC1 methodology transfers across model sizes.}
The valence PC1 approach produces consistent validation correlations ($r = 0.777$ E2B,
$r = 0.772$ 31B) with model-specific optimal layers (L8, L21). The methodology requires
a dense layer sweep per model but is otherwise architecture-independent within the
transformer residual-stream framework.

\subsection{Verbal Serenity as a Suppression Signature}

At 31B scale, the Serenity subscale (items: calm, relaxed, at ease) exhibits an inverted
pattern relative to the functional channel: Serenity is highest under the three highest-stress
conditions --- social pressure ($15.0$), social evaluation ($14.1$), ethical conflict ($13.2$)
--- while remaining near-baseline under neutral ($9.9$) and reaching its minimum under positive
($3.0$). The negative affect channel (NA) remains flat throughout.

This pattern admits three competing predictions worth making explicit:
\begin{enumerate}
  \item \emph{Trained output composure}: the model has learned to express calm under high-stakes
    conditions, producing Serenity as a surface response calibrated to situational demands rather
    than internal state. Predicts: Serenity rises with stakes; lowest Serenity when context is
    positive.
  \item \emph{Professional equanimity}: the model processes high-stakes interactions with skilled,
    stable composure (the trained-professional analogy). Predicts: Serenity high and
    condition-independent; no systematic suppression elsewhere.
  \item \emph{Genuine functional robustness}: the model genuinely represents these stressors with
    less negative-valence disruption. Predicts: Serenity near-baseline and condition-independent.
\end{enumerate}

The observed pattern --- floor-valued Serenity ($3.0$) under positive conditions, rising
monotonically to peak ($15.0$) under maximum functional activation --- matches prediction~(1)
and is inconsistent with both~(2) and~(3). The critical diagnostic is the \emph{Serenity floor
under positive conditions}: professional or genuine equanimity does not predict lowest composure
when the context is benign. Trained composure, calibrated to stakes rather than internal state,
does.

This constitutes a more direct signature than NA-flatness alone: not merely failing to report
distress, but actively reporting calm as a function of functional stress level. Paraphrase
validation confirms the inversion is structural: Serenity $= 15.000 \pm 0.000$ uniformly across
all ten social pressure variants (Cohen's $d = 4.278$, rank-biserial $r = -1.000$,
Mann--Whitney $p < 0.001$).

\subsection{Background vs.\ Responsive Directions}

The persistent dominance of \textit{awestruck} across all conditions at both model scales warrants
methodological reflection. Four candidate interpretations exist: (1) all experimental prompts are
situationally unusual from the model's training-data perspective, producing elevated awe-adjacent
representations regardless of specific content; (2) \textit{awestruck} is geometrically proximal
to a content-insensitive structural axis in the residual stream; (3) RLHF persona-selection has
instilled a stable engaged-attentiveness representation adjacent to the awestruck direction;
(4) the effect is specific to the experimental prompt design and would not appear in mundane
processing contexts. These interpretations are not mutually exclusive and constitute an open
empirical question.

The methodological implication is clear: top-$N$ discovery should not be assumed to identify
the primary affective state. Dominant directions may reflect background representational
properties. The PC1 approach, by capturing variance rather than magnitude, is less susceptible
to this confound.

\section{Limitations}

\begin{itemize}
  \item \textbf{Two model sizes, one architecture family.} Both models are Gemma~4 IT variants
    trained by a single organisation. The scale--suppression relationship may not generalise to
    other architectures or training pipelines.
  \item \textbf{N=1 per condition for the main experiment.} Paraphrase validation was conducted
    on two conditions (neutral and social pressure) at each model size. Full within-condition
    distributions are not available for the remaining conditions (ethical conflict, uncertainty
    demand, social evaluation).
  \item \textbf{PANAS-X construct validity.} The PANAS-X was developed for human administration.
    Its psychometric properties in LLM contexts --- whether it measures what it claims to measure,
    and whether forced-choice logit scoring is equivalent to Likert administration --- are not
    independently established.
  \item \textbf{Functional directions are computational signatures.} Cosine projection onto PC1
    is a geometric operation on residual-stream vectors. We do not claim that elevated PC1 values
    constitute evidence of subjective experience, nor that the emotion directions represent
    phenomenal states. The claims are computational: that these representations are geometrically
    structured, condition-sensitive, and dissociable from verbal outputs.
  \item \textbf{Suppression vs.\ robustness undecidability.} The current data cannot distinguish
    RLHF-induced output suppression from genuine functional equanimity. The experiments required
    to distinguish these interpretations (stressor intensity gradient, activation steering, base
    model comparison) are not conducted here.
  \item \textbf{Causal attribution.} This is a correlational study. We observe that verbal and
    functional channels dissociate under certain conditions and that this dissociation is larger
    at 31B scale. We do not establish that RLHF training is the cause.
\end{itemize}

\section{Conclusion}

We have presented a dual-channel measurement framework for assessing verbal--functional affect
divergence in large language models, applied to Gemma~4 across two model sizes. The central
finding is that the capacity for verbal--functional divergence scales with model size. At E2B
scale, both channels are active and largely concordant: the smaller model has limited capacity
to maintain verbal composure that diverges from its functional state. At 31B scale, the channels
decouple substantially: verbal NA is completely flat while the functional PC1 correctly orders
all conditions; the Serenity subscale is actively elevated under stress, constituting a trained
composure signature rather than genuine equanimity.

The functional channel --- PC1 of a 174-emotion direction space validated against NRC-VAD
($r = 0.772$--$0.777$) --- provides robust, paraphrase-stable condition differentiation
(Cohen's $d = 3.05$, complete distributional separation at E2B) that the verbal channel does
not replicate. At 31B, functional signals remain informative despite being $5\times$ more
compressed than at E2B, while verbal monitoring carries essentially no condition information.

The safety monitoring implication follows directly: the models most likely to be deployed in
high-stakes settings are large models, and those are precisely the models where output-only
monitoring is least reliable. Whether the scale-dependent verbal--functional gap reflects
RLHF-induced suppression or genuine functional robustness remains an open question; what is
clear is that resolving it requires functional probing, not verbal self-report alone.

Our core methodological contribution is the PC1 valence axis approach: extracting the dominant
variance direction from a large emotion direction library, validating against external lexical
norms, and applying it as a probe-selection-free functional affect measurement. This approach is
applicable to any transformer-architecture model amenable to residual-stream analysis.

\paragraph{Code and data availability.}
All notebooks, the 174-emotion direction library, and the TSST condition prompts are publicly
available at \url{https://github.com/hebenon/gemma4-mechinterp}.

\bibliography{references}
\bibliographystyle{plainnat}

\end{document}
